\section{Worked correction of an invalid rule application}
\label{sec:latex-style-rule-correction}

The following tempting use of universal introduction is invalid:
\[
  \Rule[invalid \(\forall I\)]
  {P(x)\derives P(x)}
  {P(x)\derives\forall xP(x)}.
\]
The variable \(x\) occurs free in the open assumption \(P(x)\), contrary to the
side condition \(x\notin\FV(\Gamma)\).  The failure is semantic as well: on a
two-element domain, let \(P\) hold only of the element assigned to \(x\).
Then \(P(x)\) is satisfied under that assignment while
\(\forall xP(x)\) is not, so \(\set{P(x)}\notentails\forall xP(x)\):
the semantic counterpart of the conclusion judgment fails.

One corrected derivation begins with a genuinely universal premise:
\[
  \Rule[\(\forall I\)]
  {\Rule[\(\forall E\)]
    {\forall yP(y)\derives\forall yP(y)}
    {\forall yP(y)\derives P(x)}}
  {\forall yP(y)\derives\forall xP(x)}
  \side{\(x\notin\FV(\set{\forall yP(y)})\)}.
\]
Universal elimination supplies the arbitrary instance \(P(x)\); now \(x\) is
absent from the only open assumption, so universal introduction is licensed.
The correction changes the premise rather than hiding the violated side
condition.

